\documentclass[../../main.tex]{subfiles}
\begin{document}
% SOURCE p=166; printed=151; object=ch05-normed-spaces-title; region=middle-lower
\section{赋范向量空间}\label{section:ch05-normed-spaces}
% SOURCE p=166; printed=151; object=ch05-vector-space-conventions; region=lower
令 $K$ 表示 $\mathbb R$ 或 $\mathbb C$，$\mathcal X$ 为 $K$ 上的向量空间。$\mathcal X$ 的零元直接记为 $0$，依靠上下文与标量 $0\in K$ 区别。“子空间”始终指向量子空间。
\begin{definition}[生成的一维子空间与子空间之和]\label{definition:ch05-subspace-notation}
若 $x\in\mathcal X$，用 $Kx$ 表示由 $x$ 张成的一维子空间。若 $\mathcal M,\mathcal N$ 是 $\mathcal X$ 的子空间，则 $\mathcal M+\mathcal N$ 表示子空间 $\{x+y:x\in\mathcal M,\ y\in\mathcal N\}$。
\end{definition}
% SUPPLEMENT: id=v14-ch05-subspace-details; parent=ch05-vector-space-conventions; kind=source-correction+well-definedness+motivation+hypotheses; origin=independent-derivation
\begin{remark}
\textbf{零向量例外与条件说明（重构新增）。}准确地说，$Kx=\{\lambda x:\lambda\in K\}$；当 $x\ne0$ 时它才是一维，$x=0$ 时为零维子空间。原句的“一维”须带上这个例外。两个子空间的和含 $0+0$，且
\begin{align*}
\alpha(x_1+y_1)+\beta(x_2+y_2)=(\alpha x_1+\beta x_2)+(\alpha y_1+\beta y_2)\in\mathcal M+\mathcal N,
\end{align*}
故确为子空间。它容许用来自两个方向族的向量相加；并不要求分解唯一，唯一性另需交集只有零向量。
\end{remark}
% END SUPPLEMENT: id=v14-ch05-subspace-details
% SOURCE p=166; printed=151; object=ch05-seminorm-definition; region=bottom
\begin{definition}[半范数]\label{definition:ch05-seminorm}
$\mathcal X$ 上的\textbf{半范数}是函数 $x\mapsto\|x\|$，从 $\mathcal X$ 映到 $[0,\infty)$，并满足：
\begin{enumerate}
\item\label{item:ch05-seminorm-triangle} 对所有 $x,y\in\mathcal X$，$\|x+y\|\le\|x\|+\|y\|$（\textbf{三角不等式}）。
\item\label{item:ch05-seminorm-homogeneity} 对所有 $x\in\mathcal X$ 和 $\lambda\in K$，$\|\lambda x\|=|\lambda|\|x\|$。
\end{enumerate}
\end{definition}
% SOURCE p=167; printed=152; object=ch05-norm-definition; region=top
第二条显然蕴含 $\|0\|=0$。
\begin{definition}[范数与赋范向量空间]\label{definition:ch05-norm}
若一个半范数满足 $\|x\|=0$ 仅当 $x=0$ 时成立，就称它为\textbf{范数}。带有范数的向量空间称为\textbf{赋范向量空间}，也称\textbf{赋范线性空间}。
\end{definition}
% SOURCE p=167; printed=152; object=ch05-norm-metric; region=upper
\begin{proposition}\label{proposition:ch05-norm-metric}
若 $\mathcal X$ 是赋范向量空间，则 $\rho(x,y)=\|x-y\|$ 是 $\mathcal X$ 上的度量。
\end{proposition}
\begin{proof}
这是因为
\begin{align*}
\|x-z\|&\le\|x-y\|+\|y-z\|,\\
\|x-y\|&=\|(-1)(x-y)\|=\|y-x\|.
\end{align*}
\end{proof}
% SOURCE p=167; printed=152; object=ch05-norm-topology-equivalence; region=upper-middle
\begin{definition}[范数拓扑与等价范数]\label{definition:ch05-equivalent-norms}
上述度量定义的拓扑称为 $\mathcal X$ 上的\textbf{范数拓扑}。$\mathcal X$ 上的两个范数 $\|\cdot\|_1,\|\cdot\|_2$ 称为\textbf{等价}，若存在 $C_1,C_2>0$ 使
\begin{align*}
C_1\|x\|_1\le\|x\|_2\le C_2\|x\|_1\qquad(x\in\mathcal X).
\end{align*}
\end{definition}
% SOURCE p=167; printed=152; object=ch05-equivalent-norm-consequences; region=middle
等价范数给出等价的度量，因此定义相同的拓扑和相同的 Cauchy 序列。
% SUPPLEMENT: id=v14-ch05-norm-topology-details; parent=ch05-equivalent-norm-consequences; kind=proof-detail+well-definedness+motivation+hypotheses; origin=independent-derivation
\begin{proof}
\textbf{度量及等价性的细节补写（重构新增）。}齐次性取 $\lambda=0$ 给 $\|0\|=0$，取 $\lambda=-1$ 给对称性。非负性来自值域，范数的分离条件给 $\rho(x,y)=0\Longleftrightarrow x=y$；三角不等式作用于 $(x-y)+(y-z)$，得到剩余度量公理。

若两个范数满足定义中的双边估计，则对任意 $r>0$，
\begin{align*}
B_1(x,r/C_2)\subseteq B_2(x,r),\qquad B_2(x,C_1r)\subseteq B_1(x,r).
\end{align*}
因此两种开球各能在另一种任意开球中找到同心开球，拓扑相同。把双边估计应用到 $x_n-x_m$，分别以 $\varepsilon/C_2$ 和 $C_1\varepsilon$ 控制尾部，就得 Cauchy 序列相同。故两种范数下的完备性也等价。
\end{proof}
\begin{remark}
\textbf{动机与逐条件说明（重构新增）。}半范数测量向量大小，齐次性描述缩放，三角不等式控制相加；有限非负取值使这些估计有意义。允许非零向量大小为零时，只得到伪度量，例如 $\mathbb R^2$ 上 $p(x_1,x_2)=|x_1|$ 不能区分第二坐标。范数补上分离条件后才得到 Hausdorff 度量拓扑。等价范数的两个\emph{全局}正常数分别防止一种大小相对于另一种任意缩小或放大，所以不仅开集相同，Cauchy 性也相同；一般仅拓扑等价的度量没有后一保证。
\end{remark}
% END SUPPLEMENT: id=v14-ch05-norm-topology-details
% SOURCE p=167; printed=152; object=ch05-banach-definition; region=middle
\begin{definition}[Banach 空间]\label{definition:ch05-banach-space}
关于范数度量完备的赋范向量空间称为\textbf{Banach 空间}。
\end{definition}
% SOURCE p=167; printed=152; object=ch05-completion-assertion; region=middle
每个赋范向量空间都能作为稠密子空间嵌入某个 Banach 空间。一种方法是模仿利用 Cauchy 序列从 $\mathbb Q$ 构造 $\mathbb R$ 的过程；我们将在\cref{section:ch05-linear-functionals}给出一种更简单的方法。下面给出赋范空间完备性的一个有用判别。
% SOURCE p=167; printed=152; object=ch05-series-definition; region=middle
\begin{definition}[向量级数的收敛与绝对收敛]\label{definition:ch05-series-convergence}
若 $\{x_n\}$ 是 $\mathcal X$ 中的一列，称级数 $\sum_{n=1}^{\infty}x_n$ \textbf{收敛}到 $x$，若 $\sum_{n=1}^Nx_n\to x$ 当 $N\to\infty$；若 $\sum_{n=1}^{\infty}\|x_n\|<\infty$，则称级数\textbf{绝对收敛}。
\end{definition}
% SOURCE p=167; printed=152; object=ch05-banach-series-criterion; number=5.1; region=middle-lower
\begin{theorem}\label{theorem:ch05-banach-series-criterion}
赋范向量空间 $\mathcal X$ 完备，当且仅当 $\mathcal X$ 中每个绝对收敛级数都收敛。
\end{theorem}
% SOURCE p=167; printed=152; object=ch05-banach-series-proof; region=lower
\begin{proof}
若 $\mathcal X$ 完备且 $\sum_{n=1}^{\infty}\|x_n\|<\infty$，令 $S_N=\sum_{n=1}^Nx_n$。对 $N>M$，
\begin{align*}
\|S_N-S_M\|\le\sum_{n=M+1}^N\|x_n\|\longrightarrow0\quad(M,N\to\infty),
\end{align*}
所以 $\{S_N\}$ 为 Cauchy 列，因而收敛。

反过来，设每个绝对收敛级数都收敛，并设 $\{x_n\}$ 为 Cauchy 列。可选 $n_1<n_2<\cdots$，使 $m,n\ge n_j$ 时 $\|x_n-x_m\|<2^{-j}$。令 $y_1=x_{n_1}$，对 $j>1$ 令 $y_j=x_{n_j}-x_{n_{j-1}}$。则 $\sum_{j=1}^ky_j=x_{n_k}$，且
\begin{align*}
\sum_{j=1}^{\infty}\|y_j\|\le\|y_1\|+\sum_{j=1}^{\infty}2^{-j}=\|y_1\|+1<\infty.
\end{align*}
所以 $\lim_kx_{n_k}=\sum_{j=1}^{\infty}y_j$ 存在。由于 $\{x_n\}$ 为 Cauchy 列，容易验证它与子列 $\{x_{n_k}\}$ 收敛到同一极限。
\end{proof}
% SUPPLEMENT: id=v14-ch05-series-criterion-details; parent=ch05-banach-series-criterion; kind=proof-detail+motivation+hypotheses; origin=independent-derivation
\begin{proof}
\textbf{子列选择与末步补写（重构新增）。}由 Cauchy 条件，为误差 $2^{-j}$ 选一个尾部阈值 $N_j$，递归令 $n_j\ge\max(N_j,n_{j-1}+1)$，就得到原证明要求的严格递增子列（先取 $n_1\ge N_1$）。对 $j\ge2$，相邻两指标都至少为 $n_{j-1}$，所以 $\|y_j\|<2^{-(j-1)}$；这解释了原式中从 $1$ 开始的几何尾和。

记子列极限为 $x$。任给 $\varepsilon>0$，取 $N$ 使 $m,n\ge N$ 时 $\|x_n-x_m\|<\varepsilon/2$，再取 $k$ 使 $n_k\ge N$ 且 $\|x_{n_k}-x\|<\varepsilon/2$。于是每个 $n\ge N$ 都有 $\|x_n-x\|<\varepsilon$。因此全列收敛，完备性得证。
\end{proof}
\begin{remark}
\textbf{动机与逐条件说明（重构新增）。}Banach 条件保证误差能无限累积而极限仍留在空间中；级数判别把一般 Cauchy 列压缩成误差可求和的子列增量。绝对收敛仅是实数级数 $\sum\|x_n\|$ 收敛，不预先断言向量和存在。条件要求\textbf{每一个}绝对收敛级数：在有限支集序列空间 $c_{00}$ 的一致范数下，$\sum_{n\ge1}2^{-n}e_n$ 绝对收敛，而逐坐标唯一可能的极限具有无限支集，不属于 $c_{00}$，说明完备性不可直接删除。
\end{remark}
% END SUPPLEMENT: id=v14-ch05-series-criterion-details
% SOURCE p=167,168; printed=152,153; object=ch05-banach-examples; region=bottom-top
\begin{example}\label{example:ch05-banach-function-spaces}
我们已经见过一些 Banach 空间。首先，若 $X$ 为拓扑空间，则 $B(X)$ 和 $BC(X)$ 在一致范数 $\|f\|_u=\sup_{x\in X}|f(x)|$ 下是 Banach 空间。其次，若 $(X,\mathscr M,\mu)$ 为测度空间，则 $L^1(\mu)$ 在 $L^1$ 范数 $\|f\|_1=\int|f|\,\mathrm d\mu$ 下是 Banach 空间。（若把 $L^1(\mu)$ 看成单个函数的集合，$\|\cdot\|_1$ 只是半范数；把几乎处处相等的函数等同后，它才成为范数。）
\end{example}
% SOURCE p=168; printed=153; object=ch05-l1-completeness-proof; region=upper
\begin{proof}
$L^1(\mu)$ 的完备性由\cref{theorem:ch02-absolutely-integrable-series,theorem:ch05-banach-series-criterion}推出。事实上，若 $\sum_{n=1}^{\infty}\|f_n\|_1<\infty$，则\cref{theorem:ch02-absolutely-integrable-series}表明 $f=\sum_{n=1}^{\infty}f_n$ 几乎处处存在，且
\begin{align*}
\int\left|f-\sum_{n=1}^Nf_n\right|\,\mathrm d\mu
\le\sum_{n=N+1}^{\infty}\int|f_n|\,\mathrm d\mu\longrightarrow0\quad(N\to\infty).
\end{align*}
\end{proof}
% SUPPLEMENT: id=v14-ch05-banach-examples-details; parent=ch05-banach-examples; kind=proof-detail+motivation+hypotheses; origin=independent-derivation
\begin{proof}
\textbf{两个例子的完备性细节（重构新增）。}对 $B(X)$ 中的一致 Cauchy 列 $f_n$，每个 $x$ 处的标量列 Cauchy，故有极限 $f(x)$。给定 $\varepsilon>0$，取 $N$ 使 $n,m\ge N$ 时 $|f_n(x)-f_m(x)|<\varepsilon$ 对一切 $x$ 成立。固定 $n\ge N$ 让 $m\to\infty$，得 $\|f_n-f\|_u\le\varepsilon$。又 $|f(x)|\le\|f_N\|_u+1$，只需先对误差 $1$ 取一次 $N$，故 $f$ 有界。于是 $B(X)$ 完备。若各 $f_n$ 连续，对给定 $x_0$ 和误差 $\varepsilon$，选 $n$ 使 $\|f_n-f\|_u<\varepsilon/3$，再用 $f_n$ 在 $x_0$ 的连续性控制中间一项，得 $f$ 连续；所以 $BC(X)$ 也完备。

对 $L^1$，零范数当且仅当函数几乎处处为零，所以在等价类上正定；齐次性和三角不等式由积分性质得到，且不依赖代表元。对原证的级数，非负函数 $G=\sum_n|f_n|$ 由单调收敛定理满足 $\int G=\sum_n\|f_n\|_1<\infty$，因而 $G<\infty$ 几乎处处。在可测零集 $\{G=\infty\}$ 上把 $f$ 定为零，就得到一个可测代表，且 $|f|\le G$ 几乎处处，故 $f\in L^1$。原尾和估计便证明范数收敛，而不仅是逐点收敛。
\end{proof}
\begin{remark}
\textbf{条件与动机（重构新增）。}$B(X)$ 的完备性只使用标量域完备；$BC(X)$ 再用一致极限保连续。$L^1$ 的等同约定消除了测度看不见的零集差异，绝对可积尾和则把误差控制变成可求和的实数估计。这里不要求测度有限或 $\sigma$ 有限。若 $X=\varnothing$，约定唯一空函数的范数为零，两个有界函数空间均为零空间。
\end{remark}
% END SUPPLEMENT: id=v14-ch05-banach-examples-details
% SOURCE p=168; printed=153; object=ch05-more-banach-examples; region=upper
更多例子见\cref{exercise:ch05-measure-banach,exercise:ch05-ck-banach,exercise:ch05-sobolev-banach,exercise:ch05-holder-banach}和后面的各节。
% SOURCE p=168; printed=153; object=ch05-product-norm-definition; region=upper-middle
\begin{definition}[乘积范数]\label{definition:ch05-product-norm}
若 $\mathcal X,\mathcal Y$ 是赋范向量空间，则给 $\mathcal X\times\mathcal Y$ 配备\textbf{乘积范数}
\begin{align*}
\|(x,y)\|=\max(\|x\|,\|y\|)
\end{align*}
便使它成为赋范向量空间。这里 $\|x\|$ 指 $\mathcal X$ 的范数，$\|y\|$ 指 $\mathcal Y$ 的范数。
\end{definition}
% SOURCE p=168; printed=153; object=ch05-other-product-norms; region=middle
有时也使用与它等价的其他范数，例如 $\|(x,y)\|=\|x\|+\|y\|$ 或 $\|(x,y)\|=(\|x\|^2+\|y\|^2)^{1/2}$。
% SUPPLEMENT: id=v14-ch05-product-norm-details; parent=ch05-product-norm-definition; kind=well-definedness+motivation+hypotheses; origin=independent-derivation
\begin{proof}
\textbf{良定义与等价性补写（重构新增）。}三种表达式均非负、绝对齐次，而且仅在两个坐标都为零时为零。最大值型和求和型的三角不等式直接由各坐标的三角不等式得到。平方和型则先逐坐标估计，再用非负数 $a,b,c,d$ 的不等式
\begin{align*}
\sqrt{(a+c)^2+(b+d)^2}\le\sqrt{a^2+b^2}+\sqrt{c^2+d^2}.
\end{align*}
将两边平方，这等价于 $ac+bd\le\sqrt{(a^2+b^2)(c^2+d^2)}$，而后二式由 $(ad-bc)^2\ge0$ 得到。记三种范数为 $p_\infty,p_1,p_2$，则
\begin{align*}
p_\infty\le p_2\le p_1\le2p_\infty.
\end{align*}
故它们等价。最大值范数的开球正好为两个同半径开球的积；一般两个不同半径的开球积含半径为二者最小值的同心积球，所以该拓扑就是乘积拓扑。
\end{proof}
\begin{remark}
\textbf{动机与条件（重构新增）。}在有限个坐标上同时控制误差，可用最大值、总和或平方和，所得常数比较不会损失拓扑信息。这里是有限乘积；无限个坐标的最大值可能无限，且一致控制与乘积拓扑通常不同。两个因子须在同一个标量域上，才有这里的逐坐标向量运算。
\end{remark}
% END SUPPLEMENT: id=v14-ch05-product-norm-details
% SOURCE p=168; printed=153; object=ch05-quotient-space-definition; region=middle
一个相关的构造是商空间。
\begin{definition}[商空间与商范数]\label{definition:ch05-quotient-norm}
若 $\mathcal M$ 是向量空间 $\mathcal X$ 的向量子空间，它定义等价关系 $x\sim y$ 当且仅当 $x-y\in\mathcal M$。$x$ 的等价类记作 $x+\mathcal M$；所有等价类的集合，即\textbf{商空间}，记作 $\mathcal X/\mathcal M$。它是向量空间，运算为
\begin{align*}
(x+\mathcal M)+(y+\mathcal M)&=(x+y)+\mathcal M,\\
\lambda(x+\mathcal M)&=(\lambda x)+\mathcal M.
\end{align*}
若 $\mathcal X$ 赋范且 $\mathcal M$ 闭，则 $\mathcal X/\mathcal M$ 从 $\mathcal X$ 继承一个范数，称为\textbf{商范数}：
\begin{align*}
\|x+\mathcal M\|=\inf_{y\in\mathcal M}\|x+y\|.
\end{align*}
\end{definition}
% SOURCE p=168; printed=153; object=ch05-quotient-exercise-reference; region=middle-lower
更详细的讨论见\cref{exercise:ch05-quotient-norm}。
% SUPPLEMENT: id=v14-ch05-quotient-algebra-details; parent=ch05-quotient-space-definition; kind=well-definedness+motivation+hypotheses; origin=independent-derivation
\begin{proof}
\textbf{商向量空间的良定义（重构新增）。}子空间含零、对取负和相加封闭，分别给关系的自反、对称和传递性。若 $x'=x+m$、$y'=y+n$，其中 $m,n\in\mathcal M$，则 $(x'+y')-(x+y)=m+n\in\mathcal M$，且 $\lambda x'-\lambda x=\lambda m\in\mathcal M$，所以两种运算不依赖代表。向量空间公理从 $\mathcal X$ 下降，例如商零元是 $\mathcal M$，负元为 $-x+\mathcal M$。范数性质和闭性作用在\cref{exercise:ch05-quotient-norm}中逐步验证；当 $\mathcal M=\mathcal X$ 时，商只有零元，范数直接为零。
\end{proof}
\begin{remark}
\textbf{动机与条件（重构新增）。}商空间把沿 $\mathcal M$ 方向的差异视为相同，商范数便在全部代表中寻找最小可接近的大小，所以应取下确界而不要求最小值能达到。子空间条件保证运算相容；闭性保证离 $\mathcal M$ 距离为零的向量已经属于它。若子空间不闭，所得一般只有半范数，见\cref{exercise:ch05-nonclosed-quotient}。
\end{remark}
% END SUPPLEMENT: id=v14-ch05-quotient-algebra-details
% SOURCE p=168; printed=153; object=ch05-bounded-linear-map-definition; region=lower
\begin{definition}[有界线性映射]\label{definition:ch05-bounded-linear-map}
两个赋范向量空间之间的线性映射 $T:\mathcal X\to\mathcal Y$ 称为\textbf{有界}，若存在 $C\ge0$ 使
\begin{align*}
\|Tx\|\le C\|x\|\qquad(x\in\mathcal X).
\end{align*}
\end{definition}
% SOURCE p=168; printed=153; object=ch05-boundedness-distinction; region=lower
这与集合上函数的通常有界性不同，后者要求对所有 $x$ 有 $\|Tx\|\le C$。显然，非零线性映射不可能满足后一个条件，因为对所有标量 $\lambda$ 有 $T(\lambda x)=\lambda Tx$。这里的定义意味着 $T$ 在 $\mathcal X$ 的有界子集上有界。
% SOURCE p=168; printed=153; object=ch05-linear-continuity-equivalence; number=5.2; region=lower
\begin{proposition}\label{proposition:ch05-linear-continuity}
若 $\mathcal X,\mathcal Y$ 为赋范向量空间，$T:\mathcal X\to\mathcal Y$ 为线性映射，则以下条件等价：
\begin{enumerate}
\item\label{item:ch05-linear-continuous} $T$ 连续。
\item\label{item:ch05-linear-continuous-zero} $T$ 在 $0$ 连续。
\item\label{item:ch05-linear-bounded} $T$ 有界。
\end{enumerate}
\end{proposition}
% SOURCE p=168,169; printed=153,154; object=ch05-linear-continuity-proof; region=bottom-top
\begin{proof}
\subcref{proposition:ch05-linear-continuity}{item:ch05-linear-continuous}显然蕴含\subcref{proposition:ch05-linear-continuity}{item:ch05-linear-continuous-zero}。若 $T$ 在 $0\in\mathcal X$ 连续，存在 $0$ 的邻域 $U$ 使 $T(U)\subseteq\{y\in\mathcal Y:\|y\|\le1\}$；而 $U$ 包含某个以 $0$ 为中心的球 $B=\{x\in\mathcal X:\|x\|\le\delta\}$。因此 $\|x\|\le\delta$ 时 $\|Tx\|\le1$。由于 $T$ 与标量乘法交换，$\|x\|\le a$ 时有 $\|Tx\|\le a\delta^{-1}$，即 $\|Tx\|\le\delta^{-1}\|x\|$。这证明\subcref{proposition:ch05-linear-continuity}{item:ch05-linear-continuous-zero}蕴含\subcref{proposition:ch05-linear-continuity}{item:ch05-linear-bounded}。

最后，若对一切 $x$ 有 $\|Tx\|\le C\|x\|$，则只要 $\|x_1-x_2\|\le C^{-1}\varepsilon$，便有 $\|Tx_1-Tx_2\|=\|T(x_1-x_2)\|\le\varepsilon$，所以 $T$ 连续。
\end{proof}
% SUPPLEMENT: id=v14-ch05-bounded-continuity-details; parent=ch05-linear-continuity-equivalence; kind=proof-detail+source-clarification+motivation+hypotheses; origin=independent-derivation
\begin{proof}
\textbf{缩放、零常数及有界集的细节（重构新增）。}原证明中的 $\delta$ 取正数；开邻域先含某个开球，再把半径缩半，就含上述闭球。对 $x\ne0$，向量 $z=\delta x/\|x\|$ 在该闭球内，因此 $\|Tx\|=(\|x\|/\delta)\|Tz\|\le\|x\|/\delta$；$x=0$ 由线性性直接成立。

末步若 $C=0$，则 $T=0$，它直接连续，无须取 $C^{-1}$。对任意 $C\ge0$，也可统一选 $\eta=\varepsilon/(C+1)$，当 $\|x_1-x_2\|<\eta$ 时像之差小于 $\varepsilon$。这个估计与基点无关，实际上证明一致连续性。

若 $E$ 为有界集，即 $\|x\|\le R$ 对 $x\in E$，则像中范数至多 $CR$。反过来，若线性映射把每个有界集映为有界集，单位闭球的像有某界 $M$；把非零 $x$ 除以其范数，得 $\|Tx\|\le M\|x\|$，故是这里的有界算子。若 $Tx_0\ne0$，则 $\|T(nx_0)\|=n\|Tx_0\|$ 无界，验证原段关于全域有界的区别。
\end{proof}
\begin{remark}
\textbf{动机与逐条件说明（重构新增）。}线性性把任一点附近的变化化为零点附近的变化，又把小球估计缩放成全空间的比例估计，因此一点连续就足够。一般非线性连续映射没有这个性质，例如实线上的 $x\mapsto x^2$ 连续，却不满足 $|x^2|\le C|x|$ 的全局估计。结论无需空间完备；范数提供可缩放的小球与三角估计。
\end{remark}
% END SUPPLEMENT: id=v14-ch05-bounded-continuity-details
% SOURCE p=169; printed=154; object=ch05-operator-norm-definition; number=5.3; region=upper-middle
\begin{definition}[有界算子空间与算子范数]\label{definition:ch05-operator-norm}
以 $L(\mathcal X,\mathcal Y)$ 表示 $\mathcal X$ 到 $\mathcal Y$ 的全体有界线性映射所成空间。容易验证它是向量空间，且以下函数 $T\mapsto\|T\|$ 是其上的范数，称为\textbf{算子范数}：
\begin{align}\label{align:ch05-operator-norm}
\|T\|&=\sup\{\|Tx\|:\|x\|=1\}\notag\\
&=\sup\left\{\frac{\|Tx\|}{\|x\|}:x\ne0\right\}\notag\\
&=\inf\{C:\|Tx\|\le C\|x\|\text{ 对所有 }x\}.
\end{align}
\end{definition}
% SOURCE p=169; printed=154; object=ch05-operator-norm-convention; region=middle
见\cref{exercise:ch05-operator-norm}。除非另有说明，$L(\mathcal X,\mathcal Y)$ 总配备这一范数。
% SUPPLEMENT: id=v14-ch05-operator-zero-convention; parent=ch05-operator-norm-definition; kind=source-clarification+motivation+hypotheses; origin=independent-derivation
\begin{remark}
\textbf{零空间约定与动机（重构新增）。}若 $\mathcal X=\{0\}$，原式前两个上确界的集合为空，第三式也须限制 $C\ge0$ 才有通常意义。可统一定义 $\|T\|=\sup_{\|x\|\le1}\|Tx\|=\inf\{C\ge0:\|Tx\|\le C\|x\|\ \forall x\}$；零空间上唯一算子的范数为零。当 $\mathcal X\ne\{0\}$ 时，这与原三式完全一致。所有等式和范数公理将在\cref{exercise:ch05-operator-norm}完整证明。单位球是统一的“输入预算”，其像的最大可能大小给全域最优比例界；有界性保证该数有限，线性性允许用单位向量恢复全部向量。
\end{remark}
% END SUPPLEMENT: id=v14-ch05-operator-zero-convention
% SOURCE p=169; printed=154; object=ch05-operator-space-complete; number=5.4; region=middle
\begin{proposition}\label{proposition:ch05-operator-space-complete}
若 $\mathcal Y$ 完备，则 $L(\mathcal X,\mathcal Y)$ 也完备。
\end{proposition}
% SOURCE p=169; printed=154; object=ch05-operator-completeness-proof; region=middle
\begin{proof}
设 $\{T_n\}$ 为 $L(\mathcal X,\mathcal Y)$ 中的 Cauchy 列。对 $x\in\mathcal X$，$\{T_nx\}$ 在 $\mathcal Y$ 中为 Cauchy 列，因为 $\|T_nx-T_mx\|\le\|T_n-T_m\|\|x\|$。定义 $T:\mathcal X\to\mathcal Y$ 为 $Tx=\lim_nT_nx$。留给读者验证（\cref{exercise:ch05-operator-completeness}）：$T\in L(\mathcal X,\mathcal Y)$，事实上 $\|T\|=\lim_n\|T_n\|$，并且 $\|T_n-T\|\to0$。
\end{proof}
% SUPPLEMENT: id=v14-ch05-operator-complete-details; parent=ch05-operator-space-complete; kind=omitted-proof+motivation+hypotheses; origin=independent-derivation
\begin{proof}
\textbf{原文留空部分的完整补写（重构新增）。}Cauchy 列 $T_n$ 范数有界：选 $N$ 使 $n\ge N$ 时 $\|T_n-T_N\|\le1$，再把前面有限个范数并入一个上界 $M$。给定 $x,y,\alpha,\beta$，有
\begin{align*}
\|T_n(\alpha x+\beta y)-\alpha Tx-\beta Ty\|
\le |\alpha|\|T_nx-Tx\|+|\beta|\|T_ny-Ty\|\longrightarrow0.
\end{align*}
左侧第一项的极限也为 $T(\alpha x+\beta y)$，极限唯一性给 $T$ 线性。范数的反三角不等式 $|\|u\|-\|v\||\le\|u-v\|$ 由三角不等式直接推出，故取极限得 $\|Tx\|\le M\|x\|$，即 $T$ 有界。

给定 $\varepsilon>0$，选 $N$ 使 $n,m\ge N$ 时 $\|T_n-T_m\|<\varepsilon$。固定 $n\ge N$ 和任意 $x$，令 $m\to\infty$，得 $\|(T_n-T)x\|\le\varepsilon\|x\|$。于是 $\|T_n-T\|\le\varepsilon$，证明算子范数收敛。再由 $|\|T_n\|-\|T\||\le\|T_n-T\|$，得原文所说的范数极限。
\end{proof}
\begin{remark}
\textbf{动机与条件（重构新增）。}先逐个输入取极限，再用同一个 Cauchy 阈值恢复对所有单位输入的一致估计；若只有逐点收敛而没有算子范数 Cauchy 条件，这一步不能照用。只要求值域 $\mathcal Y$ 完备，是因为所有实际取极限的 $T_nx$ 都在那里；定义域无需完备。若希望对任意定义域都成立，值域完备也不可省去：取 $\mathcal X=K$，算子 $T_y(\lambda)=\lambda y$ 将 $\mathcal Y$ 等距识别为 $L(K,\mathcal Y)$。
\end{remark}
% END SUPPLEMENT: id=v14-ch05-operator-complete-details
% SOURCE p=169; printed=154; object=ch05-operator-composition; region=middle-lower
\begin{proposition}\label{proposition:ch05-operator-composition}
若 $T\in L(\mathcal X,\mathcal Y)$，$S\in L(\mathcal Y,\mathcal Z)$，则 $ST\in L(\mathcal X,\mathcal Z)$，且 $\|ST\|\le\|S\|\|T\|$。特别地，$L(\mathcal X,\mathcal X)$ 是一个代数。
\end{proposition}
\begin{proof}
因为 $\|STx\|\le\|S\|\|Tx\|\le\|S\|\|T\|\|x\|$。
\end{proof}
% SOURCE p=169; printed=154; object=ch05-banach-algebra-definition; region=lower
\begin{definition}[Banach 代数]\label{definition:ch05-banach-algebra}
一个既是代数又是 Banach 空间的空间，若任意乘积的范数至多为各因子范数之积，则称为\textbf{Banach 代数}。
\end{definition}
\begin{example}\label{example:ch05-banach-algebras}
若 $\mathcal X$ 完备，则 $L(\mathcal X,\mathcal X)$ 是 Banach 代数。另一个例子是拓扑空间 $X$ 上的 $BC(X)$，取逐点乘法和一致范数。
\end{example}
% SOURCE p=169; printed=154; object=ch05-isomorphism-isometry-definition; region=lower
\begin{definition}[可逆算子、同构与等距映射]\label{definition:ch05-isomorphism-isometry}
若 $T\in L(\mathcal X,\mathcal Y)$ 为双射，且 $T^{-1}$ 有界，则称 $T$ \textbf{可逆}，或称为\textbf{同构}。换句话说，存在 $C>0$ 使 $\|Tx\|\ge C\|x\|$。若对所有 $x\in\mathcal X$ 有 $\|Tx\|=\|x\|$，则称 $T$ 为\textbf{等距映射}。
\end{definition}
% SOURCE p=169; printed=154; object=ch05-isometry-range; region=lower
等距映射是单射，但未必满射；不过，它是到其值域的同构。
% SUPPLEMENT: id=v14-ch05-operator-algebra-isometry-details; parent=ch05-operator-composition; kind=proof-detail+motivation+hypotheses; origin=independent-derivation
\begin{proof}
\textbf{代数与逆映射断言的细节（重构新增）。}复合线性映射仍线性，原范数估计给其有界。复合满足结合律，对算子加法满足左右分配律，并与标量乘法相容，所以 $L(\mathcal X,\mathcal X)$ 是代数；其完备性由\cref{proposition:ch05-operator-space-complete}得到。对 $BC(X)$，连续有界函数之积仍连续有界，且逐点 $|f(x)g(x)|\le\|f\|_u\|g\|_u$，故 $\|fg\|_u\le\|f\|_u\|g\|_u$，配合已证完备性即为 Banach 代数。

在 $T$ 已知双射时，其逆线性。若逆有界，选正常数 $M$ 使 $\|T^{-1}y\|\le M\|y\|$，取 $y=Tx$ 得 $\|Tx\|\ge M^{-1}\|x\|$。反过来，由正下界得到 $\|T^{-1}y\|\le C^{-1}\|y\|$。注意正下界本身并不保证满射，双射假设必须保留。

若 $T$ 等距，则 $Tx=Ty$ 蕴含 $\|x-y\|=\|T(x-y)\|=0$，故单射。它到 $T(\mathcal X)$ 的映射满射，逆保持范数，因此是同构。例如 $\mathbb R\to\mathbb R^2$，$x\mapsto(x,0)$，在 Euclid 范数下等距，却不满射到整个 $\mathbb R^2$。
\end{proof}
\begin{remark}
\textbf{动机与逐条件说明（重构新增）。}算子乘法是作用的先后复合，次乘性保证连续控制累积误差；Banach 条件还容许在代数中对绝对收敛的算子级数取和。同构要求正反两向都连续，故不能仅用代数上的双射代替；等距是更强的精确保长条件，已经保证到像的逆有界，但不负责填满目标空间。这里范数下界的“换句话说”始终在原有双射条件之下理解。
\end{remark}
% END SUPPLEMENT: id=v14-ch05-operator-algebra-isometry-details
% SOURCE p=169; printed=154; object=ch05-exercises-title; region=lower
\Needspace{12\baselineskip}
\subsection*{习题}
\subfile{section-01-exercises.tex}
\end{document}
